%!TEX root = TheoretischePhysikC.tex

\section{Vorlesung: Zentralpotential}
Man betrachte Kugelsymmetrisches Potential in drei Dimensionen:
\begin{equation}
	V(\mathbf{r}) = V(r) \quad\text{mit $r = |\mathbf{r}|$.}
\end{equation}
Das Ziel ist es, die Kinematik und Dynamik eines Teilchens, das sich in dem Potential bewegt, zu lösen. Mittels Schrödingergleichung ist
\begin{equation}
	\widehat{H}|\Psi\rangle = E|\Psi\rangle
\end{equation}
zu lösen. Diese Eigenvektorgleichung liefert uns die Eigenwerte und Eigenvektoren $\widehat{H}$, die man verwenden kann, um den Unitäroperator $U_t$ zu berechnen. Eigenwerte sind hier mehr als Technik: gegenüber den atomaren Anregungsenergien ist Raumtemperatur winzig, das Elektron sitzt fast immer im Grundzustand.

\subsection{Separation der Variablen}
Die Positionsbasis $x,y,z$ wäre nicht falsch, nur schlecht gewählt: man richtet die Darstellung nach der Symmetrie.

Die Matrixelemente des Impulsoperators in Kugelkoordinaten lauten $\forall |\Psi\rangle \in \mathcal{H}$:
\begin{equation}
	\begin{split}
		\langle r,\theta,\phi|\frac{\widehat{\mathbf{p}}^2}{2m}|\Psi\rangle &= -\frac{\hbar^2}{2m}\nabla^2 \Psi(r,\theta,\phi) \\
		&= -\frac{\hbar^2}{2m} \left[\left(\frac{\partial}{\partial r} + \frac1r\right)^2 + \frac{1}{r^2\sin(\theta)}\left(\sin(\theta)\frac{\partial}{\partial \theta}\right)^2 +  \frac{1}{r^2\sin^2(\theta)}\frac{\partial^2}{\partial \phi^2}\right]\Psi(r,\theta,\phi) \\
		&= \langle r,\theta,\phi|\frac{1}{2m}\left(\widehat{p_r}^2 + \frac{\widehat{\mathbf{L}}^2}{r^2}\right)|\Psi\rangle,
	\end{split}
\end{equation}
wobei $\Psi(r,\theta,\phi) \equiv \langle r,\theta,\phi|\Psi\rangle$ und $\langle r,\theta,\phi|\widehat{p_r}|\Psi\rangle = -i\hbar \left(\frac{\partial}{\partial r} + \frac1r\right)\Psi(r,\theta,\phi)$ den radialen Impulsoperator darstellt. Der Zusatzterm $\frac1r$ ist wesentlich -- der radiale Impuls ist \emph{nicht} $-i\hbar\frac{\partial}{\partial r}$.

Die dritte Zeile ist der Angelpunkt des Kapitels: der Winkelanteil der kinetischen Energie \emph{ist} $\widehat{\mathbf{L}}^2/(2mr^2)$.

Daher werden die Eigenwerte und Eigenvektoren $\widehat{H}$ durch
\begin{equation}\label{eq:centralpotsse}
	\langle r,\theta,\phi|\widehat{H}|\Psi\rangle = \left(\frac{\widehat{p_r}^2}{2m} + \frac{\widehat{\mathbf{L}}^2}{2mr^2} + V(\widehat{\mathbf{r}})\right)\Psi(r,\theta,\phi) = E\Psi(r,\theta,\phi)
\end{equation}
gegeben. Achten Sie auf den Hut: im Matrixelement wird aus dem Operator $V(\widehat{\mathbf{r}})$ die Zahl $V(r)$.

\begin{greybox}[Symmetrie und Erhaltungsgröße]
	Eine Drehung führt das System in sich selbst über, und zu jeder Symmetrie gehört eine Erhaltungsgröße -- der Satz von Noether, den Sie aus Theoretische Physik B als Ladungserhaltung kennen. Die Drehimpulsoperatoren \emph{erzeugen} die Rotationen; weil $\widehat{H}$ rotationsinvariant ist, muss der Drehimpuls erhalten sein.
\end{greybox}

Es gilt:
\begin{equation}
	[\widehat{H}, \widehat{\mathbf{L}}^2] = [\widehat{H}, \widehat{L}_z] = [\widehat{\mathbf{L}}^2,\widehat{L}_z] = 0.
\end{equation}
Diese Beobachtung hat eine wichtige Auswirkung: Es gibt gemeinsame Eigenvektoren für die Operatoren $\widehat{H}$, $\widehat{L}_z$ und $\widehat{\mathbf{L}}^2$ (dies ist eine allgemeine Implikation, die für alle Mengen von kommutierenden Operatoren gilt).

Ziel ist es, die gemeinsame Eigenfunktionen von $\widehat{H}$ von $\widehat{H}$, $\widehat{L}_z$ und $\widehat{\mathbf{L}}^2$ zu finden. Wenn auf diese Weise eine eindeutige Wahl von Basiszuständen des Hilbertraums erreicht wird, spricht man von einem vollständigen Satz kommutierender Observablen.

\begin{intuition}
	Normalerweise wird ein Problem schwerer, wenn man mehr fordert; hier wird es leichter. Kommutierende Observablen sind gleichzeitig scharf messbar und verhalten sich, grob gesagt, wie unabhängige Teilsysteme -- deshalb zerfällt die Lösung. Ein Separationsansatz ist kein Glücksfall, sondern die Rechenform einer Symmetrie.
\end{intuition}

Um die Eigenvektoren und Eigenwerte von $\widehat{H}$ zu berechnen, wenden wir einen Separationsansatz an:
\begin{equation}
	\Psi(r,\theta,\phi) = \text{Y}_{\ell}^m(\theta,\phi) R(r) 
\end{equation}
mit
\begin{equation}
	\widehat{\mathbf{L}}^2\text{Y}_{\ell}^m(\theta,\phi) = \hbar^2 \ell(\ell+1)\text{Y}_{\ell}^m(\theta,\phi).
\end{equation}
Wegen Normierung der $\text{Y}_{\ell}^m(\theta,\phi)$ gilt
\begin{equation}
	\iiint |\Psi(\mathbf{r})|^2\,dxdydz = \int_0^{\infty} |R(r)|^2 r^2\, dr.
\end{equation}
Setzen wir den Separationsansatz in (\ref{eq:centralpotsse}) ein, erhalten wir
\begin{equation}
	\left(\frac{\widehat{p_r}^2}{2m} + \frac{\hbar^2\ell(\ell+1)}{2mr^2} + V(r)\right)R(r) = E R(r).
\end{equation}
Der Operator $\widehat{p}_r^2$ hat die Matrixelemente 
\begin{equation}
	\begin{split}
		\langle r,\theta,\phi|\widehat{p}_r^2|\Psi\rangle &= -\hbar^2\left(\frac{\partial}{\partial r} + \frac{1}{r}\right)\left(\frac{\partial}{\partial r} + \frac{1}{r}\right)\Psi(r,\theta,\phi) \\
		&= -\hbar^2\left(\frac{1}{r}\frac{\partial}{\partial r} r\right)\left(\frac{1}{r}\frac{\partial}{\partial r} r\right)\Psi(r,\theta,\phi) \\ 
		&= -\hbar^2 \frac{1}{r}\frac{\partial^2}{\partial r^2}\left(r \Psi(r,\theta,\phi)\right).
	\end{split}
\end{equation}
Mit Hilfe von der Definition $u(r)\equiv rR(r)$ können wir die radiale Wirkung wie folgt erhalten
\begin{equation}
	\left(-\frac{\hbar^2}{2m}\frac{d^2}{dr^2} + \frac{\hbar^2\ell(\ell+1)}{2mr^2} + V(r) - E\right)u(r) = 0.
\end{equation}
Das ist ein eindimensionales Problem: $r$ spielt die Rolle der Position, der Zusatzterm $\hbar^2\ell(\ell+1)/(2mr^2)$ ist die Zentrifugalbarriere -- wie bei den Planetenbahnen.

Man kann zeigen, dass nur Lösungen mit $u(0)=0$ physikalisch sinnvoll sind, da sonst entweder das Normquadrat 
\begin{equation}
	\int_0^\infty |u(r)|^2\, dr = \infty
\end{equation}
beträgt, oder den Erwartungswert der kinetischen Energie wegen $\langle \widehat{p_r}^2\rangle = \infty$ unendlich werden lässt. Daher ist die Radialgleichung äquivalent zu einem eindimensionalen Problem auf $x>0$ mit Randbedingung bei $x=0$. Festgelegt ist $V(r)$ damit noch nicht; die wichtigste Wahl, das Coulombpotential, folgt in der nächsten Vorlesung.

